% GENERATED by build_sample.py from the week files. Do not edit by hand.
\bpart{SAMPLE}{W13 sample questions, solved (copied from the weeks)}{17}
{\small The ten questions of the W13 revision lecture, in order, each with the model answer from its week. Marks are not published; indicative only. The week page has the method, traps and variations.\par}
\begin{multicols}{2}
\hd{Sample Q1 \textmd{\small(W2, p.\,\pageref{p:W2})}}
\begin{examq}[W13 Q1: FK of an RRR arm by ETS \pg{4--5 of W13} \mk{marks not published; indicative}]
Write down the transformation matrix relating the tool frame $\{T\}$ to the station frame $\{S\}$ for the RRR manipulator below. Indicate intermediate frames. (Link data: $L_1$, $L_2$ vertical, $L_3$ upper arm, $L_4$ forearm; $L=(0.6,0.5,0.4,0.3)$ m for the numbers.)
\slidefig[0.78\linewidth]{w13-lec}{4}{29 86 326 205}{W13 Q1 arm with the lecturer's frames and ETS}
\solution
\textbf{Frames.} $\{S\}$ on the base, $z_S$ up along joint 1. $\{1\}$: $\{S\}$ turned $\theta_1$ about $z_S$. $\{2\}$: shoulder, height $L_1+L_2$ above $\{S\}$, with $y_2$ along the axis of joint 2. $\{3\}$: elbow, $L_3$ along the upper arm. $\{T\}$: tip, $L_4$ further; at $\theta=\mathbf 0$ $x_T$ points along the forearm and $y_T$ up.\\
\textbf{General rule:} ${}^S\T_T={}^S\T_1\,{}^1\T_2\,{}^2\T_3\,{}^3\T_T$ (the lecturer's ETS, W13 p.4). The arrows of joints 2 and 3 lift the arm, which is \emph{negative} about $+y$, hence the minus signs:
\[{}^S\T_T=\Rm_z(\theta_1)\;t_z(L_1{+}L_2)\;\Rm_y(-\theta_2)\;t_x(L_3)\;\Rm_y(-\theta_3)\;t_x(L_4)\;\Rm_x(90^\circ)\]
\textbf{Intermediate frames} ($c_i=\cos\theta_i$, $s_i=\sin\theta_i$):
\[{}^S\T_1=\begin{bmatrix}c_1&-s_1&0&0\\s_1&c_1&0&0\\0&0&1&0\\0&0&0&1\end{bmatrix},\qquad{}^1\T_2=\begin{bmatrix}c_2&0&-s_2&0\\0&1&0&0\\s_2&0&c_2&L_1{+}L_2\\0&0&0&1\end{bmatrix}\]
\[{}^2\T_3=\begin{bmatrix}c_3&0&-s_3&L_3\\0&1&0&0\\s_3&0&c_3&0\\0&0&0&1\end{bmatrix},\qquad{}^3\T_T=\begin{bmatrix}1&0&0&L_4\\0&0&-1&0\\0&1&0&0\\0&0&0&1\end{bmatrix}\]
\end{examq}
\begin{examq}[W13 Q1 (cont.): multiplication, result, check]
\textbf{Multiply step by step} ($c_{23}=\cos(\theta_2+\theta_3)$, $s_{23}=\sin(\theta_2+\theta_3)$, using $c_2c_3-s_2s_3=c_{23}$, $s_2c_3+c_2s_3=s_{23}$):
\[{}^S\T_2={}^S\T_1{}^1\T_2=\begin{bmatrix}c_1c_2&-s_1&-c_1s_2&0\\s_1c_2&c_1&-s_1s_2&0\\s_2&0&c_2&L_1{+}L_2\\0&0&0&1\end{bmatrix}\]
\[{}^S\T_3={}^S\T_2{}^2\T_3=\begin{bmatrix}c_1c_{23}&-s_1&-c_1s_{23}&c_1c_2L_3\\s_1c_{23}&c_1&-s_1s_{23}&s_1c_2L_3\\s_{23}&0&c_{23}&L_1{+}L_2{+}s_2L_3\\0&0&0&1\end{bmatrix}\]
Last factor: $\Rm_x(90^\circ)$ sends column 2 of the rotation to $\ +$column 3 and column 3 to $-$column 2, and $t_x(L_4)$ adds $L_4\times$ column 1 to the position. With $r=L_3c_2+L_4c_{23}$:
\[\boxed{{}^S\T_T=\begin{bmatrix}c_1c_{23}&-c_1s_{23}&s_1&c_1\,r\\s_1c_{23}&-s_1s_{23}&-c_1&s_1\,r\\s_{23}&c_{23}&0&L_1{+}L_2{+}L_3s_2{+}L_4s_{23}\\0&0&0&1\end{bmatrix}}\]
\textbf{Check at $\theta=\mathbf 0$:} tip $(L_3{+}L_4,\,0,\,L_1{+}L_2)$ and $x_T=x_S$, $y_T=z_S$, $z_T=-y_S$, as drawn. At $\theta_2=90^\circ$, $\theta_3=0$: tip $(0,0,L_1{+}L_2{+}L_3{+}L_4)$, arm straight up (as $\Rm_y(-\theta_2)$ predicts).\\
\textbf{Numbers} $\theta=(30^\circ,20^\circ,40^\circ)$: frame $\{2\}$ at $(0,0,1.1)$, frame $\{3\}$ at $(0.3255,\,0.1879,\,1.2368)$ m,
\[{}^S\T_T=\begin{bmatrix}0.4330&-0.7500&0.5000&0.4554\\0.2500&-0.4330&-0.8660&0.2629\\0.8660&0.5000&0&1.4966\\0&0&0&1\end{bmatrix}\ \chk\]
\textbf{Marks lost by:} $\Rm_y(+\theta_2)$ (arm dips when $\theta_2$ grows), dropping the final $\Rm_x(90^\circ)$ (position right, orientation wrong), giving only the product with no intermediate frames when the question asks for them.
\end{examq}
\hd{Sample Q2 \textmd{\small(W3, p.\,\pageref{p:W3})}}
\begin{examq}[W13 Q2 \mk{marks not published; $\sim$1 mark per step}]

\emph{Sketch the workspace of the manipulator (RRR arm of Q1: rotation $\theta_1$ about $z_S$, then two joints about horizontal axes, links $L_3,L_4$ above a shoulder at height $L_1+L_2$). Assume no limit on joint angles, ignore collisions.}
\slidefig[0.8\linewidth]{w13-lec}{6}{38 97 96 135}{the question and the lecturer's ``donut'': a ring in the side view}
\solution
\textbf{Frames.} $\{S\}$ at the base, $z_S$ up. Shoulder point $O=(0,0,L_1+L_2)$.\\
\textbf{Reasoning.} Joints 2 and 3 are a planar RR arm (links $L_3,L_4$) in a vertical plane through $z_S$, pivoting about $O$. By the cosine rule the tip is at distance $\rho=\sqrt{L_3^2+L_4^2+2L_3L_4\cos\theta_3}$ from $O$; with $\theta_3$ free, $\rho$ takes every value from $|L_3-L_4|$ (folded) to $L_3+L_4$ (stretched). With $\theta_2$ free, the tip sweeps a full annulus in that plane. Joint 1 spins the plane about $z_S$, which passes through $O$.\\
\textbf{Result.} The workspace is a \textbf{hollow sphere (spherical shell) centred on $O$}, outer radius $L_3+L_4$, inner radius $|L_3-L_4|$; $L_1,L_2$ only lift the centre. For $L_3=0.4$, $L_4=0.3$ m: $0.1\le\rho\le0.7$ m about $O$. If $L_3=L_4$ the hole closes: a solid ball. The dexterous workspace is empty (a position fixes the tool direction up to the branch).
\begin{center}
\begin{tikzpicture}[scale=1.0]
\def\R{1.75}\def\r{0.25}
\begin{scope}
\fill[g2,even odd rule] (0,2.6) circle (\R) (0,2.6) circle (\r);
\draw[line width=0.6pt] (0,2.6) circle (\R); \draw[line width=0.6pt] (0,2.6) circle (\r);
\draw[line width=0.8pt] (-1,0)--(1,0); \draw[line width=0.8pt] (0,0)--(0,2.6);
\fill (0,2.6) circle (0.04); \node[font=\scriptsize,anchor=west] at (0.05,2.45) {$O$};
\draw[<->,line width=0.4pt] (0.2,0.05)--(0.2,2.55) node[midway,right,font=\scriptsize]{$L_1{+}L_2$};
\draw[->,line width=0.4pt] (0,2.6)--(\R,2.6) node[pos=0.55,above,font=\tiny]{$L_3{+}L_4$};
\node[font=\scriptsize] at (0,-0.35) {side view ($x_S z_S$ plane)};
\end{scope}
\begin{scope}[xshift=5.0cm,yshift=2.6cm]
\fill[g2] (0,0) circle (\R); \draw[line width=0.6pt] (0,0) circle (\R); \draw[line width=0.5pt,dashed] (0,0) circle (\r);
\fill (0,0) circle (0.04); \draw[->,line width=0.4pt] (0,0)--(\R,0) node[pos=0.6,above,font=\tiny]{$L_3{+}L_4$};
\node[font=\scriptsize] at (0,-2.95) {top view (disc)};
\end{scope}
\end{tikzpicture}
\end{center}
\drawncap{side view: ring about $O$ (hole radius $|L_3-L_4|$, to scale; height of $O$ not to scale). Top view: full disc of radius $L_3+L_4$ (hole not visible from above: those points are reached at other heights).}
\textbf{Check} at a simple pose: $\theta_2=\theta_3=0$, arm horizontal, tip at $(L_3+L_4,0,L_1+L_2)$, on the outer sphere ✓; folded $\theta_3=\pi$: $\rho=|L_3-L_4|$ ✓ (500 random poses give $\rho\in[0.1,0.7]$). \chk
\end{examq}
\hd{Sample Q3 \textmd{\small(W3, p.\,\pageref{p:W3})}}
\begin{examq}[W13 Q3 \mk{marks not published}]

\emph{Derive the inverse kinematics of the RRR arm of Q1. You are given the position of the end effector $(x_T,y_T,z_T)$.}
\solution
\textbf{FK (from Q1).} ${}^S\mathbf p_T=(c_1r,\ s_1r,\ L_1+L_2+L_3s_2+L_4s_{23})$, $r=L_3c_2+L_4c_{23}$; $c_{23}=\cos(\theta_2+\theta_3)$.
\par \textbf{Top-down view} gives $\theta_1$: $\theta_1=\atantwo(y_T,x_T)$, $r=+\sqrt{x_T^2+y_T^2}$ (arm faces the point). Branch 2: $\theta_1+\pi$ with $r=-\sqrt{x_T^2+y_T^2}$ (arm reaches back over the base axis). atan2, not the $\tan^{-1}(y/x)$ of the slide.\\
\par \textbf{Side-on view}: planar RR in the plane $(r,z')$, $z'=z_T-L_1-L_2$:
\[r=L_3c_2+L_4c_{23},\qquad z'=L_3s_2+L_4s_{23}\]
Square and add (the cross terms give $2L_3L_4(c_2c_{23}+s_2s_{23})=2L_3L_4\cos\theta_3$):
\begin{gather*}r^2+z'^2=L_3^2+L_4^2+2L_3L_4\cos\theta_3\ \Rightarrow\ \cos\theta_3=\frac{r^2+z'^2-L_3^2-L_4^2}{2L_3L_4},\\ \theta_3=\atantwo\big(\pm\sqrt{1-\cos^2\theta_3},\cos\theta_3\big)\end{gather*}
Expand $c_{23},s_{23}$ and write $r=(L_3+L_4c_3)c_2-(L_4s_3)s_2$, $z'=(L_4s_3)c_2+(L_3+L_4c_3)s_2$; the point $(r,z')$ is $(L_3+L_4c_3,\ L_4s_3)$ rotated by $\theta_2$:
\[\theta_2=\atantwo(z',r)-\atantwo(L_4\sin\theta_3,\ L_3+L_4\cos\theta_3)\]
\par \textbf{Number and existence.} $2$ (shoulder) $\times\,2$ (elbow) $=4$ for an interior point. No solution if $|\cos\theta_3|>1$, i.e.\ the point lies outside the shell of Q2 (outside $|L_3-L_4|\le\rho\le L_3+L_4$ about the shoulder). On the shell surface the two elbow branches merge. If $x_T=y_T=0$, $\theta_1$ is arbitrary.\\
\par \textbf{Numbers.} $L=(0.6,0.5,0.4,0.3)$ m, $\mathbf p_T=(0.4554,0.2629,1.4966)$ m. $\theta_1=\atantwo(0.2629,0.4554)=30^\circ$, $r=0.5259$, $z'=1.4966-1.1=0.3966$; $\cos\theta_3=\dfrac{0.2766+0.1573-0.16-0.09}{0.24}=0.7660$, $\theta_3=\pm40^\circ$; $\sigma=\atantwo(0.3966,0.5259)=37.024^\circ$, $\mu=\pm17.024^\circ$. Solutions $(\theta_1,\theta_2,\theta_3)$ in degrees:
\[(30,\ 20,\ 40),\quad(30,\ 54.047,\ -40),\quad(-150,\ 125.953,\ 40),\quad(-150,\ 160,\ -40)\]
Each reproduces $\mathbf p_T$ by FK; the first is the Q1 pose, so the check closes. Reach test: a target $(0.5,0.5,1.8)$ is $0.995$ m from the shoulder $>0.7$: no solution. \chk\par
\textbf{Sanity:} $\theta_2=\theta_3=0$ gives $r=L_3+L_4$, $z=L_1+L_2$ (stretched, on the outer sphere) ✓.
\end{examq}
\hd{Sample Q4 \textmd{\small(W7, p.\,\pageref{p:W7})}}
\begin{examq}[W13 Q4 (a): derive the equations of motion, 3D RRR arm \mk{marks not published; indicative, ~1 per step}]
\slidefig[0.5\linewidth]{w13-lec}{8}{48 108 610 200}{W13 p.\,8: RRR arm; the lecturer's red dots mark $m_1$ (shoulder), $m_2$ (elbow), $m_3$ (tool). $\theta_1$ about $z_S$, $\theta_2,\theta_3$ in the vertical arm plane}
``Mass, moments of inertia and positions of centre of mass are given. Derive the equations of motion.''
\solution
\textbf{Set-up (state it).} Frame $\{S\}$: $z_S$ up, gravity along $-z_S$. $\theta_1$ turns the whole arm about $z_S$. The arm plane is vertical, at azimuth $\theta_1$; shoulder at height $H=L_1+L_2$ on the axis. $\theta_2$: elevation of the upper arm (length $L_3$) from the horizontal; $\theta_3$: forearm (length $L_4$) relative to the upper arm, so the forearm's absolute elevation is $\theta_2+\theta_3$. (If the exam measures angles from the vertical, swap $c\leftrightarrow s$ in heights.) Point masses first (the red dots): $m_2$ at the elbow, $m_3$ at the tool; $m_1$ sits on the axis, so it only adds a constant to $P$ and nothing to the equations. Use $\rho$ = horizontal distance from $z_S$ and $h$ = height.\\
\textbf{Positions.} $\rho_2=L_3c_2$, $h_2=H+L_3s_2$; $\rho_3=L_3c_2+L_4c_{23}$, $h_3=H+L_3s_2+L_4s_{23}$.\\
\textbf{Velocities.} A point at $(\rho,h)$ on a plane spinning at $\dot\theta_1$: $v^2=\dot\rho^2+\dot h^2+\rho^2\dot\theta_1^2$. $\dot\rho_2^2+\dot h_2^2=L_3^2\dot\theta_2^2$. For mass 3: $(\dot\rho_3,\dot h_3)=L_3\dot\theta_2(-s_2,c_2)+L_4(\dot\theta_2+\dot\theta_3)(-s_{23},c_{23})$, so $\dot\rho_3^2+\dot h_3^2=L_3^2\dot\theta_2^2+L_4^2(\dot\theta_2+\dot\theta_3)^2+2L_3L_4c_3\,\dot\theta_2(\dot\theta_2+\dot\theta_3)$ (using $s_2s_{23}+c_2c_{23}=c_3$).\\
\textbf{Energies.}
\begin{align*}K&=\tfrac12\big(m_2\rho_2^2+m_3\rho_3^2\big)\dot\theta_1^2+\tfrac12\big[(m_2+m_3)L_3^2+m_3L_4^2+2m_3L_3L_4c_3\big]\dot\theta_2^2\\&\quad+m_3\big(L_4^2+L_3L_4c_3\big)\dot\theta_2\dot\theta_3+\tfrac12m_3L_4^2\dot\theta_3^2\\
P&=m_2g(H+L_3s_2)+m_3g(H+L_3s_2+L_4s_{23})+\text{const},\qquad \mathcal L=K-P\end{align*}
Reading $\mathbf M$ off $K=\tfrac12\dot{\boldsymbol\theta}^\top\mathbf M\dot{\boldsymbol\theta}$: $M_{11}=m_2\rho_2^2+m_3\rho_3^2$; $M_{22}=(m_2+m_3)L_3^2+m_3L_4^2+2m_3L_3L_4c_3$; $M_{23}=m_3(L_4^2+L_3L_4c_3)$; $M_{33}=m_3L_4^2$; $M_{12}=M_{13}=0$.
\end{examq}
\begin{examq}[W13 Q4 (b): Euler--Lagrange lines \mk{cont.}]
\textbf{Joint 1.} $\partial\mathcal L/\partial\dot\theta_1=M_{11}\dot\theta_1$, $\ \frac d{dt}=M_{11}\ddot\theta_1+\dot M_{11}\dot\theta_1$; $\partial\mathcal L/\partial\theta_1=0$ (nothing depends on $\theta_1$). With $\dot\rho_2=-L_3s_2\dot\theta_2$, $\dot\rho_3=-L_3s_2\dot\theta_2-L_4s_{23}(\dot\theta_2+\dot\theta_3)$, $\dot M_{11}=2m_2\rho_2\dot\rho_2+2m_3\rho_3\dot\rho_3$:
\[Q_1=M_{11}\ddot\theta_1-2\dot\theta_1\Big[m_2\rho_2L_3s_2\dot\theta_2+m_3\rho_3\big(L_3s_2\dot\theta_2+L_4s_{23}(\dot\theta_2+\dot\theta_3)\big)\Big]\]
\textbf{Joint 2.} $\partial\mathcal L/\partial\dot\theta_2=M_{22}\dot\theta_2+M_{23}\dot\theta_3$; $\frac d{dt}=M_{22}\ddot\theta_2+M_{23}\ddot\theta_3-m_3L_3L_4s_3\dot\theta_3(2\dot\theta_2+\dot\theta_3)$ (using $\dot M_{22}=-2m_3L_3L_4s_3\dot\theta_3$, $\dot M_{23}=-m_3L_3L_4s_3\dot\theta_3$). $\partial\mathcal L/\partial\theta_2=\tfrac12\dot\theta_1^2\,\partial M_{11}/\partial\theta_2-g\big[(m_2+m_3)L_3c_2+m_3L_4c_{23}\big]$ with $\tfrac12\partial M_{11}/\partial\theta_2=-\big[m_2\rho_2L_3s_2+m_3\rho_3(L_3s_2+L_4s_{23})\big]$.
\begin{multline*}Q_2=M_{22}\ddot\theta_2+M_{23}\ddot\theta_3+\dot\theta_1^2\big[m_2\rho_2L_3s_2+m_3\rho_3(L_3s_2+L_4s_{23})\big]\\-m_3L_3L_4s_3(2\dot\theta_2\dot\theta_3+\dot\theta_3^2)+g\big[(m_2+m_3)L_3c_2+m_3L_4c_{23}\big]\end{multline*}
\textbf{Joint 3.} $\partial\mathcal L/\partial\dot\theta_3=M_{23}\dot\theta_2+M_{33}\dot\theta_3$; $\frac d{dt}=M_{23}\ddot\theta_2+M_{33}\ddot\theta_3-m_3L_3L_4s_3\dot\theta_3\dot\theta_2$. $\partial K/\partial\theta_3=-m_3\rho_3L_4s_{23}\dot\theta_1^2-m_3L_3L_4s_3(\dot\theta_2^2+\dot\theta_2\dot\theta_3)$; $\partial P/\partial\theta_3=m_3gL_4c_{23}$. The $\dot\theta_2\dot\theta_3$ terms cancel:
\[Q_3=M_{23}\ddot\theta_2+M_{33}\ddot\theta_3+m_3\rho_3L_4s_{23}\dot\theta_1^2+m_3L_3L_4s_3\dot\theta_2^2+g\,m_3L_4c_{23}\]
\end{examq}
\begin{examq}[W13 Q4 (c): $\mathbf M,\mathbf C,\mathbf g$, general links, numbers \mk{cont.}]
\[\mathbf M=\begin{bmatrix}M_{11}&0&0\\0&M_{22}&M_{23}\\0&M_{23}&M_{33}\end{bmatrix},\quad
\mathbf g=\begin{bmatrix}0\\g[(m_2+m_3)L_3c_2+m_3L_4c_{23}]\\g\,m_3L_4c_{23}\end{bmatrix}\]
\[\mathbf C=(V_1,V_2,V_3)^\top=\text{the velocity terms of }Q_1,Q_2,Q_3\]
Centrifugal: the $\dot\theta_1^2$ terms in $V_2,V_3$ and $m_3L_3L_4s_3\dot\theta_2^2$ (squares). Coriolis: $\dot\theta_1\dot\theta_2$, $\dot\theta_1\dot\theta_3$ in $V_1$; $\dot\theta_2\dot\theta_3$ in $V_2$ (products). Physical reading: spinning about $z_S$ flings the arm outwards; raising the arm while spinning changes the inertia about $z_S$, so joint 1 must push; gravity loads joints 2 and 3 only, most when the arm is horizontal.\\
\textbf{With link inertias} (slender links, COM at $l_2$ from the shoulder and $l_3$ from the elbow, $I_2,I_3$ about axes $\perp$ the link through the COM, $\approx0$ about the link axis, $I_1$ the column about $z_S$): replace $L_3\to l_2$ in link 2's terms and $L_4\to l_3$ in link 3's (keep $L_3$ where the elbow position enters $\rho_3$), $\rho_2=l_2c_2$, $\rho_3=L_3c_2+l_3c_{23}$, and
\[M_{11}=I_1+m_2\rho_2^2+m_3\rho_3^2+I_2c_2^2+I_3c_{23}^2,\quad M_{22}=m_2l_2^2+I_2+m_3(L_3^2+l_3^2+2L_3l_3c_3)+I_3\]
\[M_{23}=m_3(l_3^2+L_3l_3c_3)+I_3,\quad M_{33}=m_3l_3^2+I_3,\quad \mathbf g=\begin{bmatrix}0\\g[(m_2l_2+m_3L_3)c_2+m_3l_3c_{23}]\\g\,m_3l_3c_{23}\end{bmatrix}\]
\[V_1=-2\dot\theta_1\big[m_2\rho_2l_2s_2\dot\theta_2+m_3\rho_3(L_3s_2\dot\theta_2+l_3s_{23}(\dot\theta_2+\dot\theta_3))+I_2c_2s_2\dot\theta_2+I_3c_{23}s_{23}(\dot\theta_2+\dot\theta_3)\big]\]
\[V_2=\dot\theta_1^2\big[m_2\rho_2l_2s_2+m_3\rho_3(L_3s_2+l_3s_{23})+I_2c_2s_2+I_3c_{23}s_{23}\big]-m_3L_3l_3s_3(2\dot\theta_2\dot\theta_3+\dot\theta_3^2)\]
\[V_3=\dot\theta_1^2\big[m_3\rho_3l_3s_{23}+I_3c_{23}s_{23}\big]+m_3L_3l_3s_3\dot\theta_2^2\]
(The red-dot case is $l_2=L_3$, $l_3=L_4$, $I_2=I_3=0$.) Rotational energy added: $\tfrac12I_1\dot\theta_1^2+\tfrac12I_2(\dot\theta_2^2+c_2^2\dot\theta_1^2)+\tfrac12I_3((\dot\theta_2+\dot\theta_3)^2+c_{23}^2\dot\theta_1^2)$.\\
\textbf{Numbers} (point masses): $m_2=2$, $m_3=1$ kg, $L_3=0.6$, $L_4=0.4$ m, $\boldsymbol\theta=(20^\circ,30^\circ,-40^\circ)$, $\dot{\boldsymbol\theta}=(1,0.5,-0.5)$ rad/s, $\ddot{\boldsymbol\theta}=(0.2,1,0.5)$ rad/s$^2$. $\rho_2=0.5196$, $\rho_3=0.9135$ m, $c_3=0.7660$, $s_3=-0.6428$, $c_{23}=0.9848$, $s_{23}=-0.1736$.
\[\mathbf M=\begin{bmatrix}1.3746&0&0\\0&1.6077&0.3439\\0&0.3439&0.1600\end{bmatrix},\ \mathbf C=\begin{bmatrix}-0.5858\\0.4838\\-0.1020\end{bmatrix},\ \mathbf g=\begin{bmatrix}0\\19.157\\3.864\end{bmatrix}\]
$Q_1=1.3746(0.2)-0.5858=\boxed{-0.311}$; $Q_2=1.6077(1)+0.3439(0.5)+0.4838+19.157=\boxed{21.420}$; $Q_3=0.3439(1)+0.16(0.5)-0.1020+3.864=\boxed{4.186}$ N\,m. Hold still: $\mathbf Q=\mathbf g=(0,19.157,3.864)$; arm stretched out horizontally: $(0,21.582,3.924)$ N\,m (largest). \chk\\
\textbf{Checks.} $\mathbf M$ symmetric, positive definite; $\mathbf g$ contains no $\dot\theta$; $\mathbf C=\mathbf 0$ at rest; power $\dot{\boldsymbol\theta}^\top\mathbf Q=8.306$ W $=\dot K+\dot P$; free swing conserves $K+P$ to $10^{-9}$ J over 3 s. \chk
\end{examq}
\hd{Sample Q5 \textmd{\small(W4, p.\,\pageref{p:W4})}}
\slidefig[0.85\linewidth]{w13-lec}{9}{0 76 48 92}{W13 p.\,9: the question (K and $\phi$ given; find $\theta_1,\varphi_1,\varphi_2,\varphi_3$)}
\begin{examq}[W13 Q5: four-bar IK \mk{marks not published}]
A planar four-bar has $L_0=1.0$, $L_1=0.5$, $L_3=0.6$ m; the output link (coupler) has $a=0.2$, $b=0.25$ m. Point K on the output link is at $(0.3710,\,0.4678)$ m in {S}, and the link is oriented at $\phi=-0.600^\circ$. Determine $\theta_1,\varphi_1,\varphi_2,\varphi_3$.
\solution
Frames: {S} at the left pivot, {D} at $(L_0,0)$, angles counter-clockwise from $+x$. Write each chain's K equations: $x=L_1\cos\theta_1+a\cos(\theta_1+\varphi_2)$, $y=L_1\sin\theta_1+a\sin(\theta_1+\varphi_2)$, $\theta_1+\varphi_2=\phi$ (mod $2\pi$).\mk{setup}\\
Substitute: $L_1\cos\theta_1=x-a\cos\phi$, $L_1\sin\theta_1=y-a\sin\phi$: $\theta_1=\operatorname{atan2}(0.4699,0.1710)=\boxed{70.0^\circ}$; $\varphi_2=\phi-\theta_1=\boxed{289.4^\circ}$ ($\equiv-70.6^\circ$).\mk{left chain}\\
Right chain: $x=L_0+L_3\cos\varphi_1+b\cos(\varphi_1+\varphi_3)$ with $\varphi_1+\varphi_3=\phi+\pi$: $L_3\cos\varphi_1=x-L_0+b\cos\phi=-0.3790$, $L_3\sin\varphi_1=y+b\sin\phi=0.4652$: $\varphi_1=\operatorname{atan2}(0.4652,-0.3790)=\boxed{129.2^\circ}$; $\varphi_3=\phi+\pi-\varphi_1=\boxed{50.2^\circ}$.\mk{right chain}\\
Check: $\lVert(-0.379,0.4652)\rVert=0.600=L_3$; $\theta_1+\varphi_2=359.4^\circ=\varphi_1+\varphi_3+180^\circ$. Say: IK is easy because each leg is solved alone with atan2; one answer per leg. Using the slide's printed $\varphi_3=\varphi_1-\phi+\pi$ gives $309.8^\circ$ and K misses by 0.4 m.\mk{check}
\end{examq}
\hd{Sample Q6 \textmd{\small(W1, p.\,\pageref{p:W1})}}
\begin{examq}[W13 Q6: mobility of a planar parallel manipulator \mk{marks not published}]

\slidefig[0.55\linewidth]{w13-lec}{13}{48 113 480 157}{W13 pp.12--13: the same figure; on p.12 the lecturer numbers the links 1 ground, 2 platform, 3--4 left leg, 5--6 middle leg, 7--8 right leg}
Find the mobility of the mechanism drawn (planar, three R--P--R legs).
\solution
$M=D(N-J-1)+\sum_{i=1}^{J}F_i$ (G--K).\quad Planar, so $D=3$.\\
Links: 1 ground, 2 platform, and each prismatic joint joins a cylinder to a rod (two links per leg: 3,4; 5,6; 7,8): $N=8$.\\
Joints: per leg R, P, R: $J=9$, each $F_i=1$, $\sum F_i=9$.
\[M=3(8-9-1)+9=-6+9=\boxed{3}\qquad\text{(form 1: }3(8-1)-2\cdot9=3\text{)}\]
Three DoF: platform $x$, $y$ and rotation, driven by the three prismatic actuators. Check: $N=7$ (ground forgotten) would give $0$. \chk
\end{examq}
\hd{Sample Q7 \textmd{\small(W4, p.\,\pageref{p:W4})}}
\slidefig[0.85\linewidth]{w13-lec}{14}{0 76 67 108}{W13 p.\,14: the delta mechanism and its joint graph (B base, R, S, S, M moving platform)}
\begin{examq}[W13 Q7: delta mobility, plain vs extended \mk{marks not published}]
A delta robot has a base, three R joints, each followed by a parallelogram of two parallel bars with spherical joints at both ends of every bar, and a moving platform. (a) Find the mobility with plain G--K. (b) The platform has only 3 translational DoF. Reconcile with the extended equation.
\solution
(a) 3D, $D=6$. $N=1+3+6+1=11$. $J=3+3\times4=15$. $\sum F_i=3(1+4\times3)=39$.
\[M=6(11-15-1)+39=\boxed{9}\]
(b) Each of the 6 bars has an S joint at both ends, so it can spin about its own axis with nothing else moving: $R_M=6$. No constraint is repeated: $R_C=0$.
\[M=6(11-15-1)+39+0-6=\boxed{3}\]
Explain: plain G--K counts every DoF as moving the output; the 6 bar spins are internal, so they are subtracted. The result 3 equals the 3 motors and the 3 platform translations.
\end{examq}
\hd{Sample Q8 \textmd{\small(W5, p.\,\pageref{p:W5})}}
\begin{examq}[W13 Q8, symbolic via-point cubics \mk{marks not published}]
\slidefig{w13-lec}{18}{29 68 48 146}{W13 p.\,18: question and the lecturer's set-up (the slide's cubics are $\theta_1=a_0{+}a_1t{+}a_2t^2{+}a_3t^3$ and $\theta_2=b_0{+}b_1t{+}b_2t^2{+}b_3t^3$)}
\solution
\textbf{Assume} (state it): start and end at rest, each segment on its own clock, equal durations $t_f$; frames: joint angle $\theta$ in rad (or deg).\\
\textbf{Conditions (8 for 8 unknowns):} $\theta_1(0)=\theta_0$, $\dot\theta_1(0)=0$, $\theta_1(t_f)=\theta_v$, $\theta_2(0)=\theta_v$, $\theta_2(t_f)=\theta_g$, $\dot\theta_2(t_f)=0$, $\dot\theta_1(t_f)=\dot\theta_2(0)$, $\ddot\theta_1(t_f)=\ddot\theta_2(0)$.\\
\textbf{Step 1.} $a_0=\theta_0$, $a_1=0$, $b_0=\theta_v$. Let $w=\dot\theta_2(0)=\dot\theta_1(t_f)$.\\
\textbf{Step 2.} Segment 1 has end speeds $(0,w)$ and $\Delta=\theta_v-\theta_0$: $a_2=\dfrac{3\Delta-wt_f}{t_f^2}$, $a_3=\dfrac{-2\Delta+wt_f}{t_f^3}$. Segment 2 has $(w,0)$ and $\Delta'=\theta_g-\theta_v$: $b_1=w$, $b_2=\dfrac{3\Delta'-2wt_f}{t_f^2}$, $b_3=\dfrac{-2\Delta'+wt_f}{t_f^3}$.\\
\textbf{Step 3.} $\ddot\theta_1(t_f)=2a_2+6a_3t_f=\dfrac{-6\Delta+4wt_f}{t_f^2}$, $\ddot\theta_2(0)=2b_2=\dfrac{6\Delta'-4wt_f}{t_f^2}$. Equal: $8wt_f=6(\Delta+\Delta')\Rightarrow$ $\boxed{w=\dfrac{3(\theta_g-\theta_0)}{4t_f}}$.\\
\textbf{Step 4.} Substitute $w$:
\[\boxed{\begin{aligned}
a_2&=\tfrac{12\theta_v-3\theta_g-9\theta_0}{4t_f^2}, & a_3&=\tfrac{-8\theta_v+3\theta_g+5\theta_0}{4t_f^3},\\
b_1&=\tfrac{3(\theta_g-\theta_0)}{4t_f}, & b_2&=\tfrac{-12\theta_v+6\theta_g+6\theta_0}{4t_f^2},\quad b_3=\tfrac{8\theta_v-5\theta_g-3\theta_0}{4t_f^3}\end{aligned}}\]
\textbf{Check} with $\theta=(10,40,20)^\circ$, $t_f=1$\,s: $w=7.5^\circ$/s, $\theta_1=10+82.5t^2-52.5t^3$, $\theta_2=40+7.5t-75t^2+47.5t^3$; $\theta_1(1)=40$, $\theta_2(1)=20$, $\ddot\theta_1(1)=-150=\ddot\theta_2(0)$ \checkmark. Variant $(0,40,30)^\circ$, $t_f=1$: $w=22.5$; segment 1 $97.5t^2-57.5t^3$, segment 2 $40+22.5t-75t^2+42.5t^3$; peak speed in segment 1 $55.11^\circ$/s at $t=0.565$\,s. \chk\\
\textbf{Common slips:} one cubic for both segments; forgetting $a_1=0$; using total time instead of $t_f$ per segment; dropping the acceleration condition (then $w$ is undetermined).
\end{examq}
\hd{Sample Q9 \textmd{\small(W5, p.\,\pageref{p:W5})}}
\begin{examq}[W13 Q9, cubic for a prismatic joint \mk{marks not published}]
\slidefig{w13-lec}{19}{38 308 48 130}{W13 p.\,19 (question text)}
\solution
Prismatic joint, so $q=d$ in m. Conditions: $d(0)=0.02$, $d(0.8)=0.15$, $\dot d(0)=0$, $\dot d(0.8)=0$. $\Delta d=0.13$\,m, $t_F=0.8$\,s.\\
$a_0=0.02$\,m, $a_1=0$, $a_2=\dfrac{3(0.13)}{0.8^2}=0.609375$\,m/s$^2$, $a_3=\dfrac{-2(0.13)}{0.8^3}=-0.5078125$\,m/s$^3$.
\[\boxed{d(t)=0.02+0.609375\,t^2-0.5078125\,t^3\ \text{m}}\qquad \dot d=1.21875t-1.5234375t^2\ \text{m/s}\]
Check: $d(0.8)=0.02+0.39-0.26=0.15$ \checkmark; $\dot d(0.8)=0.975-0.975=0$ \checkmark. Peak speed $1.5\Delta d/t_F=0.24375$\,m/s at $t=0.4$\,s ($d=0.085$\,m); end accelerations $\pm1.21875$\,m/s$^2$. \chk\\
\textbf{Common slips:} $a_3$ divided by $t_F^2$; cm/m mixed; forgetting $a_0=0.02$.
\end{examq}
\hd{Sample Q10 \textmd{\small(W10, p.\,\pageref{p:W10})}}
\begin{examq}[(W13 Q10) \mk{marks not published; indicative ~1 per step}]
Consider a robot with three omnidirectional wheels, each independently driven, rollers perpendicular to each wheel's driving direction. The robot is an equilateral triangle; the distance from the centre to each vertex is $d$. Derive a kinematic expression that relates the angular velocities of the wheels $(\dot\varphi_1,\dot\varphi_2,\dot\varphi_3)$ to the velocity of the robot chassis frame $(u,v,\dot\theta)$. Is the robot holonomic? Find the wheel speeds for $u=0.2$ m/s with $d=0.10$ m, $r=0.05$ m.
\slidefig[0.5\linewidth]{w13-lec}{20}{29 151 634 119}{W13 p.\,20: the question's figure}
\solution
\textbf{Setup.} $\{C\}$ at the centre, $\hat x_C$ to the right, $\hat y_C$ up. Wheels at $W_1(0,d)$, $W_2(\tfrac{\sqrt3}2d,-\tfrac d2)$, $W_3(-\tfrac{\sqrt3}2d,-\tfrac d2)$; drive directions tangential (clockwise, as drawn): $\beta_1=0^\circ$, $\beta_2=-120^\circ$, $\beta_3=120^\circ$; $\gamma_i=0$ (rollers $\perp$ drive). \mk{2}\\
\textbf{General formula (slide \pg{36}).} $\dot\varphi_i=\begin{bmatrix}\frac1r&\frac{\tan\gamma_i}r\end{bmatrix}\mathbf R(\beta_i)^\top\begin{bmatrix}1&0&-\wXy i\\0&1&\wXx i\end{bmatrix}\wXz$; with $\gamma_i=0$: $r\dot\varphi_i=\cos\beta_i(u-\dot\theta\wXy i)+\sin\beta_i(v+\dot\theta\wXx i)$. \mk{1}\\
\textbf{Rows.} $r\dot\varphi_1=u-d\dot\theta$; $r\dot\varphi_2=-\frac12(u+\frac d2\dot\theta)-\frac{\sqrt3}2(v+\frac{\sqrt3}2d\dot\theta)=-\frac12u-\frac{\sqrt3}2v-d\dot\theta$; $r\dot\varphi_3=-\frac12u+\frac{\sqrt3}2v-d\dot\theta$. \mk{3}
\[\boxed{\begin{bmatrix}\dot\varphi_1\\\dot\varphi_2\\\dot\varphi_3\end{bmatrix}=\frac1r\begin{bmatrix}1&0&-d\\-\frac12&-\frac{\sqrt3}2&-d\\-\frac12&\frac{\sqrt3}2&-d\end{bmatrix}\begin{bmatrix}u\\v\\\dot\theta\end{bmatrix}}\]
\textbf{Inverse and holonomic.} $\det=\tfrac{3\sqrt3}2d\neq0$; $\wXz=\frac r3\begin{bmatrix}2&-1&-1\\0&-\sqrt3&\sqrt3\\-\frac1d&-\frac1d&-\frac1d\end{bmatrix}\wXp$. Swedish wheels give no no-slide row, so $u,v,\dot\theta$ are free: holonomic, $\delta_m=3$. World frame: $\wXn=\mathbf J(\theta)\wXz$. \mk{2}\\
\textbf{Numbers.} $u=0.2$: $\dot\varphi=\frac1{0.05}(0.2,-0.1,-0.1)=\boxed{(4,\,-2,\,-2)\ \text{rad/s}}$. \mk{1}\\
\textbf{Check.} Wheel 1 drives along $\hat x_C$ so it carries all of $u$; wheels 2 and 3 drive at $\pm120^\circ$ and see $-\frac12u$; a pure spin gives $-d\dot\theta/r$ on every wheel; sum of the three speeds is $0$ for pure translation. \chk
\end{examq}
\end{multicols}
